
\documentclass[12pt]{article}
\usepackage{amsmath,graphicx,geometry}
\geometry{margin=1in}

\title{Rapport Scientifique -- Estimation du Temps de Livraison}
\author{}
\date{}

\begin{document}
\maketitle

\section{Introduction}
Les plateformes de livraison doivent estimer le temps de livraison afin d'optimiser la satisfaction client et la logistique. La problématique traitée est la prédiction supervisée du temps de livraison (régression).

\section{Méthodologie}
\subsection{Pré-traitement}
Nettoyage des doublons, conversion des dates, suppression des valeurs aberrantes. Imputation par médiane pour les variables numériques et par mode pour les catégorielles. Encodage One-Hot et standardisation des variables numériques.

\subsection{EDA}
Analyse des distributions, heatmap des corrélations, identification des variables influentes (charge du système, disponibilité des livreurs, type de restaurant).

\subsection{Modélisation}
Trois modèles testés : Régression Linéaire, Random Forest, XGBoost. Validation par cross-validation et optimisation par RandomizedSearchCV.

\section{Résultats}
XGBoost est le meilleur modèle. Les métriques utilisées : MAE, RMSE et $R^2$. Les erreurs proviennent principalement des heures de pointe et de la variabilité logistique.

\section{Conclusion}
Modèle performant mais limites dues à l'absence de données géographiques et temporelles fines. Améliorations : ajout de features temporelles, géolocalisation, modèles CatBoost et LightGBM.

\end{document}
